% Options for packages loaded elsewhere
\PassOptionsToPackage{unicode}{hyperref}
\PassOptionsToPackage{hyphens}{url}
%
\documentclass[
]{article}
\usepackage{amsmath,amssymb}
\usepackage{iftex}
\ifPDFTeX
  \usepackage[T1]{fontenc}
  \usepackage[utf8]{inputenc}
  \usepackage{textcomp} % provide euro and other symbols
\else % if luatex or xetex
  \usepackage{unicode-math} % this also loads fontspec
  \defaultfontfeatures{Scale=MatchLowercase}
  \defaultfontfeatures[\rmfamily]{Ligatures=TeX,Scale=1}
\fi
% \usepackage{lmodern} (commented for build)
\ifPDFTeX\else
  % xetex/luatex font selection
\fi
% Use upquote if available, for straight quotes in verbatim environments
\IfFileExists{upquote.sty}{\usepackage{upquote}}{}
\IfFileExists{XXmicrotype.sty}{% use microtype if available
  \usepackage[expansion=false,protrusion=false]{microtype}
  \UseMicrotypeSet[protrusion]{basicmath} % disable protrusion for tt fonts
}{}
\makeatletter
\@ifundefined{KOMAClassName}{% if non-KOMA class
  \IfFileExists{parskip.sty}{%
    \usepackage{parskip}
  }{% else
    \setlength{\parindent}{0pt}
    \setlength{\parskip}{6pt plus 2pt minus 1pt}}
}{% if KOMA class
  \KOMAoptions{parskip=half}}
\makeatother
\usepackage{xcolor}
\usepackage{longtable,booktabs,array}
\usepackage{calc} % for calculating minipage widths
% Correct order of tables after \paragraph or \subparagraph
\usepackage{etoolbox}
\makeatletter
\patchcmd\longtable{\par}{\if@noskipsec\mbox{}\fi\par}{}{}
\makeatother
% Allow footnotes in longtable head/foot
\IfFileExists{footnotehyper.sty}{\usepackage{footnotehyper}}{\usepackage{footnote}}
\makesavenoteenv{longtable}
\setlength{\emergencystretch}{3em} % prevent overfull lines
\providecommand{\tightlist}{%
  \setlength{\itemsep}{0pt}\setlength{\parskip}{0pt}}
\setcounter{secnumdepth}{-\maxdimen} % remove section numbering
\ifLuaTeX
  \usepackage{selnolig}  % disable illegal ligatures
\fi
\IfFileExists{bookmark.sty}{\usepackage{bookmark}}{\usepackage{hyperref}}
\IfFileExists{xurl.sty}{\usepackage{xurl}}{} % add URL line breaks if available
\urlstyle{same}
\hypersetup{
  hidelinks,
  pdfcreator={LaTeX via pandoc}}

\author{}
\date{}

\begin{document}

\hypertarget{a-parameter-free-prediction-of-the-solar-mixing-angle-from-cubic-group-flavor-structure-tested-by-juno}{%
\section{A Parameter-Free Prediction of the Solar Mixing Angle from
Cubic-Group Flavor Structure, Tested by
JUNO}\label{a-parameter-free-prediction-of-the-solar-mixing-angle-from-cubic-group-flavor-structure-tested-by-juno}}

\textbf{Author:} Alex Maybaum\\
\textbf{Date:} April 2026\\
\textbf{Status:} DRAFT PRE-PRINT\\
\textbf{Classification:} Theoretical Physics / High Energy Physics /
Neutrino Phenomenology

\begin{center}\rule{0.5\linewidth}{0.5pt}\end{center}

\hypertarget{abstract}{%
\subsection{Abstract}\label{abstract}}

We derive the parameter-free prediction $\sin^{2}\theta_{12}$ = 1/3 $-$ 1/($4\pi^{2}$) =
0.3080 from cubic-group flavor structure, with no parameter fit to the
value of $\sin^{2}\theta_{12}$. The first JUNO measurement of reactor neutrino
oscillations reports $\sin^{2}\theta_{12}$ = 0.3092 ± 0.0087, with the post-JUNO
global fit tightening this to $\sin^{2}\theta_{12}$ = 0.3085 ± 0.0073, matching the
prediction at 0.$07\sigma$. The same construction gives
$\sin^{2}\theta_{13}$ = 4/($18\pi^{2}$) = 0.02252 and $\sin^{2}\theta_{23}$ = 1/2 + 1/($2\pi^{2}$) = 0.5507, each
consistent with NuFIT 6.0 {[}14{]} within 1.$1\sigma$, and forces the exact
sum rule 2$\sin^{2}\theta_{12}$ + $\sin^{2}\theta_{23}$ = 7/6 --- a clean discriminator from the
TM1 and TM2 patterns currently analyzed against JUNO. The angle
predictions follow from a discrete $\mu$-$\tau$ parity together with one
cubic-group structural relation between the U-even and U-odd
components of the second-order perturbation. The Cabibbo angle
$\lambda$ = 1/($\pi\sqrt{2}$) = 0.22508, set by the cubic Brillouin-zone geometry, sets
the perturbation scale and matches the observed value to 0.04\%.
JUNO's design-lifetime precision of ±0.0014 will test the prediction
at sub-percent level.

\begin{center}\rule{0.5\linewidth}{0.5pt}\end{center}

\hypertarget{introduction}{%
\subsection{1. Introduction}\label{introduction}}

\hypertarget{the-result}{%
\subsubsection{1.1 The result}\label{the-result}}

The first JUNO measurement of reactor neutrino oscillations {[}7{]} reports

\textbf{$\sin^{2}\theta_{12}$ = 0.3092 ± 0.0087}

from 59.1 days of data. The post-JUNO global fit of Capozzi \emph{et
al.} {[}8{]} tightens this to

\textbf{$\sin^{2}\theta_{12}$ = 0.3085 ± 0.0073}.

We show that taking the cubic point group O --- the rotation symmetry
group of three-dimensional space --- as the underlying flavor structure
of the lepton sector yields the prediction

\textbf{$\sin^{2}\theta_{12}$ = 1/3 $-$ 1/($4\pi^{2}$) = 0.30798}

matching the post-JUNO global fit at \textbf{0.$07\sigma$} with no parameter
fit to $\sin^{2}\theta_{12}$. The same construction predicts $\sin^{2}\theta_{13}$ and $\sin^{2}\theta_{23}$
within 1.$1\sigma$ of the NuFIT 6.0 {[}14{]} best fits and forces the exact sum rule

\textbf{2$\sin^{2}\theta_{12}$ + $\sin^{2}\theta_{23}$ = 7/6}

which distinguishes the present prediction from the TM1, TM2, and TM3
column-preservation patterns at JUNO precision. Both inputs to $\sin^{2}\theta_{12}$
are independently motivated: the zeroth-order value 1/3 is the $A_{4}$ TBM
result, and the perturbation scale $\lambda^{2}$ = 1/($2\pi^{2}$) is the squared Cabibbo
angle, fixed independently by the cubic Brillouin geometry (matching the
observed $\lambda$ to 0.04\%).

\hypertarget{where-the-prediction-comes-from}{%
\subsubsection{1.2 Where the prediction comes
from}\label{where-the-prediction-comes-from}}

The structure is constrained by two ingredients. First, the perturbation
scale $\lambda^{2}$ = 1/($2\pi^{2}$) is not a fitted parameter: it is the squared Cabibbo
angle, fixed by a single cubic Brillouin-zone distance (Sec. 2.2).
Numerically $\lambda$ = 1/($\pi\sqrt{2}$) = 0.22508, agreeing with the observed Cabibbo
value 0.22500 ± 0.00067 {[}13{]} to 0.04\%. Second, the deviations
from tribimaximal (TBM) mixing at O($\lambda^{2}$) are controlled by a discrete $\mu$-$\tau$
parity (the U-parity of $S_{4}\supset A_{4}$) together with one cubic-group
structural relation between the U-even and U-odd components of the
second-order perturbation (Sec. 3.4). The structural relation fixes the
coefficient of the 1/($4\pi^{2}$) correction; without further inputs, the angle
predictions are

\begin{itemize}
\tightlist
\item
  \textbf{$\sin^{2}\theta_{12}$ = 1/3 $-$ 1/($4\pi^{2}$) = 0.30798}
\item
  \textbf{$\sin^{2}\theta_{13}$ = 4/($18\pi^{2}$) = 0.02252}
\item
  \textbf{$\sin^{2}\theta_{23}$ = 1/2 + 1/($2\pi^{2}$) = 0.55066}
\end{itemize}

The 0.$07\sigma$ match with the post-JUNO global fit therefore tests structural
relations rather than fitted parameters.

\hypertarget{the-empirical-context}{%
\subsubsection{1.3 The empirical context}\label{the-empirical-context}}

The cubic-group structure used here is not a postulated flavor symmetry.
It is the rotation symmetry group of three-dimensional space, realized
in a lattice-based framework {[}18{]} in which the simple
cubic Bravais lattice in d=3 is shown to be the unique three-dimensional
Bravais lattice consistent with the SU(3) $\times$ SU(2) $\times$ U(1) gauge structure
and the observed Cabibbo angle (Theorem 7b). The same cubic structure
that produces the 1/($4\pi^{2}$) correction here also produces a number of
related quantitative predictions across the flavor sector, all
empirically tested:

\begin{longtable}[]{@{}
  >{\raggedright\arraybackslash}p{(\columnwidth - 6\tabcolsep) * \real{0.2500}}
  >{\raggedright\arraybackslash}p{(\columnwidth - 6\tabcolsep) * \real{0.2500}}
  >{\raggedright\arraybackslash}p{(\columnwidth - 6\tabcolsep) * \real{0.2500}}
  >{\raggedright\arraybackslash}p{(\columnwidth - 6\tabcolsep) * \real{0.2500}}@{}}
\toprule\noalign{}
\begin{minipage}[b]{\linewidth}\raggedright
Observable
\end{minipage} & \begin{minipage}[b]{\linewidth}\raggedright
Cubic-group prediction
\end{minipage} & \begin{minipage}[b]{\linewidth}\raggedright
Observed
\end{minipage} & \begin{minipage}[b]{\linewidth}\raggedright
Match
\end{minipage} \\
\midrule\noalign{}
\endhead
\bottomrule\noalign{}
\endlastfoot
Cabibbo angle $\lambda$ & 1/($\pi\sqrt{2}$) = 0.22508 & 0.22500 ± 0.00067 {[}13{]} &
0.04\% \\
Koide angle $\theta_{0}$ (charged leptons) & $C_{2}/d^{2}$ = 2/9 & 0.222204 ± 0.000010
{[}13{]} & 0.02\% \\
Solar mixing $\sin^{2}\theta_{12}$ (this work) & 1/3 $-$ 1/($4\pi^{2}$) = 0.30798 & 0.3085 ±
0.0073 {[}8{]} & 0.$07\sigma$ \\
PMNS sum rule & 2$\sin^{2}\theta_{12}$ + $\sin^{2}\theta_{23}$ = 7/6 & 1.178 ± 0.020 & 0.$6\sigma$ \\
\end{longtable}

The first two entries are derived in {[}18, §§7.1--7.2{]}; the
Cabibbo derivation is unconditional structural (the chirality structure
of the taste-changing vertex is established in the staggered fermion
literature {[}18, §7.1, citing Mason et al.{]}), and the
empirical match to 0.04\% supports this status. The Koide angle is
unconditional structural at first-principles. The third and fourth
entries are the present work, classified at the layered (Layer 1 + Layer
2(b)) level per {[}18, §8.3{]} following operator-counting
analysis of Condition 2; see Sec. 6.3 for full discussion. This
cumulative empirical track record is what distinguishes the cubic-group
structure from postulated flavor symmetries: the same parameter $\lambda^{2}$ =
1/($2\pi^{2}$) and the same projection factor $A^{2}$ = 2/3 enter multiple
observables at sub-percent precision, which is testable structure rather
than symmetry assumption.

The remainder of this paper develops the technical content. Section 2
sets up the cubic-group structure and derives the perturbation scale.
Section 3 computes the deviations from TBM at O($\lambda^{2}$) and derives the sum
rule. Section 4 presents the predictions and compares them with JUNO,
the post-JUNO global fit, and competing TM1/TM2 patterns. Section 5
covers the remaining angles. Section 6 addresses robustness. Section 7
discusses falsifiability under JUNO's projected long-term precision.
Section 8 concludes.

\hypertarget{position-in-the-broader-literature}{%
\subsubsection{1.4 Position in the broader literature}\label{position-in-the-broader-literature}}

The lattice-based framework in which the cubic-group structure is
realized {[}18{]} is one of several recent constructions in which
observer-essentiality plays a foundational role. Chandrasekaran, Longo,
Penington, and Witten {[}19{]} established that incorporating an
observer promotes the von Neumann algebra of QFT observables in a
gravitational subregion from type III to type II, with a well-defined
entropy. Maldacena {[}20{]} showed that an observer-with-clock cancels
the unphysical phase of the de Sitter sphere partition function.
Harlow, Usatyuk, and Zhao {[}21{]} argued that the closed-universe
Hilbert space dimension is determined by the observer's degrees of
freedom. Slagle and Preskill {[}22{]} constructed boundary quantum
mechanics from a classical bulk lattice model with stochastic
dynamics, demonstrating the logical possibility of QM-from-classical
at the cost of an exponentially long extra spatial dimension. The
framework underlying this paper differs from these constructions in
being deterministic at the substratum level (no stochastic dynamics,
no extra spatial dimension), and in producing the cubic-group flavor
structure used here as a structural consequence rather than an input.
The convergence with the post-2022 mainstream observer-emergent
literature on the foundational observer-essentiality move is
supportive context for the present prediction; the prediction itself
stands or falls on its empirical match with JUNO, which is the focus
of this paper.

\hypertarget{cubic-group-flavor-structure}{%
\subsection{2. Cubic-group flavor
structure}\label{cubic-group-flavor-structure}}

\hypertarget{the-cubic-group-o-and-its-aux2084-subgroup}{%
\subsubsection{2.1 The cubic group O and its $A_{4}$
subgroup}\label{the-cubic-group-o-and-its-aux2084-subgroup}}

The cubic point group O is the rotation symmetry group of
three-dimensional space, the unique finite point group with a
six-element generator set \{±$\hat{e}_{1}$, ±$\hat{e}_{2}$, ±$\hat{e}_{3}$\} --- six minimal-length
translations corresponding to forward and backward motion along three
orthogonal axes. In the lattice realization of the OI framework, this
six-element set arises from coupling-degree minimization (K = 2d = 6 for
d = 3), and the multiplicities (3, 2, 1) under the action of O on these
six directions match the structure of one Standard Model fermion
generation (quark triplet, lepton doublet, singlet). Three SM
generations are placed on the $T_{1}$ representation of the cubic group O
acting on the three Cartesian axes --- a forced consequence of
{[}18, §4.6 Theorem 7{]} (cubic decomposition
\(6 = T_1 \oplus E \oplus A_1\)), {[}18, §4.6 Theorem 7b{]} (Bravais-lattice
uniqueness --- SC is the unique 3D Bravais lattice compatible with the
SM gauge group at the commutant level and with the observed Cabibbo
angle), {[}18, §4.7 Theorem 10{]} (cubic symmetry forces three degenerate
triplet sectors), {[}18, §4.7.1.2{]} (link-carrier construction placing
matter on links with cross-charged representations admissible), and
{[}18, Theorem 15{]} (anomaly cancellation fixing hypercharges uniquely).
The detailed argument from coupling-degree minimization through the (3,
2, 1) decomposition is given in {[}18, §§4.5--4.7{]}.

The full octahedral rotation group O has order 24 and is isomorphic to
$S_{4}$. Its irreducible representations decompose as

\textbf{irreps(O) = $A_{1}\oplus A_{2}\oplus$ E $\oplus T_{1}\oplus T_{2}$}

of dimensions 1, 1, 2, 3, 3. The subgroup $A_{4}\subset$ O keeps \{$A_{1}$, $A_{2}$, E, $T_{1}$\}
but identifies $A_{2}$ with $A_{1}$, giving the standard $A_{4}$ content A $\oplus$ A' $\oplus$ A'\,'
$\oplus$ T of dimensions 1, 1, 1, 3.

We take $T_{1}$ as the generation triplet. Realize it by the three Cartesian
unit vectors \{$e_{1}$, $e_{2}$, $e_{3}$\} acted on by the 24 rotations of O. The
democratic vector

\textbf{$\hat{h}$ = (1/$\sqrt{3}$)($e_{1}$ + $e_{2}$ + $e_{3}$)}

is the only direction fixed by the $\mathbb{Z}$$_{3}$ cyclic subgroup. The plane
perpendicular to $\hat{h}$ carries the doublet E. This is the geometric origin
of TBM: the second column of U\_TBM is exactly $\hat{h}$, and the first and
third columns lie in the plane perpendicular to it, oriented by $\mu$-$\tau$
exchange.

\hypertarget{the-cabibbo-scale-from-cubic-brillouin-geometry}{%
\subsubsection{2.2 The Cabibbo scale from cubic Brillouin
geometry}\label{the-cabibbo-scale-from-cubic-brillouin-geometry}}

We identify the perturbation parameter $\lambda$ with the Cabibbo angle. In a
lattice realization of the cubic-group flavor structure, the three
generations occupy the three $T_{1}$ corners of the cubic Brillouin zone at
momenta X\_j = $\pi\cdot e$\_j (j = 1, 2, 3). The inter-generation distance in
momentum space is

\textbf{\textbar X\_i $-$ X\_j\textbar{} = $\pi\sqrt{2}$ (i $\neq$ j)}

a fixed geometric constant of the cubic lattice. The leading-order
inter-generation mixing matrix element is the continuum fermion
propagator at this momentum, \textbar M\_ij\textbar{} = \textbar S(X\_j
$-$ X\_i)\textbar{} = 1/\textbar X\_j $-$ X\_i\textbar. The
1/\textbar q\textbar{} scaling (rather than 1/$|q|^{2}$) is
structural: in the massless limit, the chirality-preserving component of
the fermion propagator at momentum q reduces to S(q) = $\gamma\cdot q/q^{2}$, and the
taste-changing vertex's Dirac trace contracts this to a magnitude of
order 1/\textbar q\textbar{} rather than the full propagator's
1/$|q|^{2}$. The full derivation, including the chirality
verification of the taste-changing vertex's spinor structure, is given
in {[}18, §7.1{]}. Identifying the matrix element with the
Cabibbo angle gives

\textbf{$\lambda$ = 1/($\pi\sqrt{2}$) = 0.22508}

matching the observed value $\lambda$\_obs = 0.22500 ± 0.00067 {[}13{]} to
0.04\%. Squaring,

\textbf{$\lambda^{2}$ = 1/($2\pi^{2}$) = 0.05066}

sets the overall scale of O($\lambda^{2}$) corrections. The full derivation of $\lambda$ =
1/($\pi\sqrt{2}$) from the lattice fermion propagator and the
chirality-preservation argument is given in {[}18, §7.1{]}.
The empirical match to 0.04\% is at the level of structural relations
rather than fitted parameters.

A second geometric quantity will appear. The angle between any
generation axis e\_i and the democratic direction $\hat{h}$ satisfies e\_i $\cdot$ $\hat{h}$ =
1/$\sqrt{3}$, so $\cos^{2}$($\angle$(e\_i, $\hat{h}$)) = 1/3 and $\sin^{2}$($\angle$(e\_i, $\hat{h}$)) = 2/3. We identify
the Wolfenstein parameter A with this quantity:

\textbf{A = $\sqrt{2/3}$ = 0.8165}

agreeing with the PDG value A = 0.826 ± 0.012 at 0.$23\sigma$. The combination
$A^{2}$ = 2/3 enters the inter-generation projection geometry, and $A^{4}$ = 4/9
appears directly in the deviation $\Delta_{13}$ below.

\hypertarget{the-pmns-matrix-in-the-altarelli-feruglio-basis}{%
\subsubsection{2.3 The PMNS matrix in the Altarelli-Feruglio
basis}\label{the-pmns-matrix-in-the-altarelli-feruglio-basis}}

We work in the Altarelli-Feruglio (AF) basis {[}12{]}. In this basis the neutrino sector realizes TBM exactly (after
diagonalization of the seesaw matrix), and all deviations from TBM come
from the charged-lepton mass matrix. The PMNS matrix is

\textbf{U\_PMNS = V\_$\ell$$^{\dagger}\cdot$ U\_TBM}

where V\_$\ell$ is the charged-lepton rotation that diagonalizes M\_$\ell$ M\_$\ell$$^{\dagger}$,
and

\[
U_{\rm TBM} = \begin{pmatrix} \sqrt{2/3} & \sqrt{1/3} & 0 \\ -\sqrt{1/6} & \sqrt{1/3} & -\sqrt{1/2} \\ -\sqrt{1/6} & \sqrt{1/3} & \sqrt{1/2} \end{pmatrix}
\]

Expanding V\_$\ell$ in $\lambda$,

\textbf{V\_$\ell$ = exp($\lambda A_{1}$ + $\lambda^{2}A_{2}$ + O($\lambda^{3}$))}

with $A_{1}$, $A_{2}$ real antisymmetric 3$\times$3 matrices. The remaining task is to
determine $A_{1}$ and $A_{2}$ from the structural ingredients of the cubic
geometry.

\hypertarget{u-parity-grading-and-the-sum-rule}{%
\subsection{3. U-parity grading and the sum
rule}\label{u-parity-grading-and-the-sum-rule}}

\hypertarget{the-u-parity-generator}{%
\subsubsection{3.1 The U-parity
generator}\label{the-u-parity-generator}}

In the AF basis, the $\mu$-$\tau$ exchange generator U $\in S_{4}$ acts on the
generation indices by exchanging the second and third:

\[
U = \begin{pmatrix} 1 & 0 & 0 \\ 0 & 0 & 1 \\ 0 & 1 & 0 \end{pmatrix}
\]

Conjugation by U defines a $\mathbb{Z}$$_{2}$ grading on antisymmetric matrices: A is
\textbf{U-even} if UAU$^{-1}$ = +A, \textbf{U-odd} if UAU$^{-1}$ = $-$A. With
(E\_\{ij\})\emph{\{kl\} = $\delta$}\{ik\}$\delta$\_\{jl\} $-\delta$\_\{il\}$\delta$\_\{jk\}, the
decomposition is:

\begin{itemize}
\tightlist
\item
  \textbf{U-even}: $E_{12}$ + $E_{13}$
\item
  \textbf{U-odd}: $E_{12}-E_{13}$, $E_{23}$
\end{itemize}

The unique U-odd combination antisymmetric in ($\mu$, $\tau$) is $E_{23}$; the unique
U-odd combination relating the first generation to the second/third is
$E_{12}-E_{13}$.

\hypertarget{determination-of-aux2081}{%
\subsubsection{3.2 Determination of $A_{1}$}\label{determination-of-aux2081}}

$A_{1}$ is fixed by three structural conditions plus the projection geometry
of §2.2.

The data require $\sin\theta_{13}$ = O($\lambda$) (since $\theta_{13}$ $\neq$ 0), forcing $A_{1}$ to contribute
to the (1,3) element of V\_$\ell$$^{\dagger}$ U\_TBM. The observation that both $\Delta_{12} \equiv$
$\sin^{2}\theta_{12}-$ 1/3 and $\Delta_{23} \equiv$ $\sin^{2}\theta_{23}-$ 1/2 are O($\lambda^{2}$) (not O($\lambda$)) further
constrains $A_{1}'s$ structure. The remaining input is geometric: each end of
an inter-generation mixing vertex projects onto the plane perpendicular
to $\hat{h}$, contributing a factor sin($\angle$(e\_i, $\hat{h}$)) = A from the off-democratic
projection (Sec. 2.2). For a single-vertex inter-generation transition,
the total projection factor is $A^{2}$ (one factor per vertex end).
Identifying this projection geometry with the leading O($\lambda$) contribution
to U\_e3 gives the structural identity

\textbf{$\sin\theta_{13}$ = $A^{2}\lambda$}

(at O($\lambda$); $A_{2}$ contributes only at O($\lambda^{2}$) and cancels from
$|U_e3|^{2}$ as shown in §3.3 below). This is the structural
input --- not a derived consequence --- that fixes $A_{1}'s$ overall
coefficient up to sign.

The result is

\textbf{$A_{1}$ = ($\sqrt{2}/3$)($E_{12}-E_{13}$)}

$A_{1}$ is purely U-odd: $UA_{1}U$$^{-1}$ = $-A_{1}$. Geometrically, it is a rotation by
angle $A^{2}$ = 2/3 around the axis ($e_{2}$ + $e_{3}$)/$\sqrt{2}$, which reproduces $\sin\theta_{13}$ =
$A^{2}\lambda$ as required.

\hypertarget{deviations-to-oux3bbuxb2}{%
\subsubsection{3.3 Deviations to O($\lambda^{2}$)}\label{deviations-to-oux3bbuxb2}}

Expanding V\_$\ell$ to second order,

\textbf{V\_$\ell$ = I + $\lambda A_{1}$ + $\lambda^{2}$($A_{2}$ + 1/2$A_{1}^{2}$) + O($\lambda^{3}$)}

and writing the most general antisymmetric $A_{2}$ as

\textbf{$A_{2}$ = $b_{12}E_{12}$ + $b_{13}E_{13}$ + $b_{23}E_{23}$}

we compute $|U_PMNS,ij|^{2}$ to O($\lambda^{2}$). The PMNS angles
satisfy $\sin^{2}\theta_{13}$ = $|U_e3|^{2}$, $\sin^{2}\theta_{12}\cos^{2}\theta_{13}$ =
$|U_e2|^{2}$, and $\sin^{2}\theta_{23}\cos^{2}\theta_{13}$ = $|U_\mu 3|^{2}$.
The deviations from TBM are:

\begin{itemize}
\tightlist
\item
  \textbf{$\Delta_{13}$ = (4/9)$\lambda^{2}$ = $A^{4}\lambda^{2}$}
\item
  \textbf{$\Delta_{12}$ = $-$(2/3)($b_{12}$ + $b_{13}$) $\lambda^{2}$}
\item
  \textbf{$\Delta_{23}$ = $b_{23} \lambda^{2}$}
\end{itemize}

Three observations follow.

\textbf{(i) $\Delta_{13}$ is fully determined by $A_{1}$ alone.} The $A_{2}$ parameters
cancel from $|U_e3|^{2}$ at O($\lambda^{2}$), so $\sin^{2}\theta_{13}$ = (4/9)$\lambda^{2}$ is
fixed without reference to $A_{2}$. Equivalently, the structural identity
$\sin\theta_{13}$ = $A^{2}\lambda$ is exact at O($\lambda$) and receives no O($\lambda^{2}$) correction.

\textbf{(ii) The combination $b_{12}-b_{13}$ does not enter any of the three
angles.} It is naturally identified with the CP-violating Dirac phase
$\delta$\_CP, since it is the only remaining free parameter at O($\lambda^{2}$) that does
not enter the angles. The present analysis leaves $\delta$\_CP undetermined.

\textbf{(iii) The remaining structure depends on two combinations:} ($b_{12}$
+ $b_{13}$) controlling $\Delta_{12}$ and $b_{23}$ controlling $\Delta_{23}$. The first is the U-even,
the second the dominant U-odd coefficient of $A_{2}$, with $b_{12}-b_{13}$ the
sub-dominant U-odd combination.

\hypertarget{sum-rule-and-absolute-normalization}{%
\subsubsection{3.4 Sum rule and absolute
normalization}\label{sum-rule-and-absolute-normalization}}

The deviations depend on two scalar combinations of $A_{2}$: ($b_{12}$ + $b_{13}$)
controlling $\Delta_{12}$, and $b_{23}$ controlling $\Delta_{23}$. We fix the normalization with
two conditions.

\textbf{Condition 1 (natural normalization).} Take $b_{23}$ = 1. This sets $A_{2}$
to enter the perturbation series at unit strength relative to $\lambda^{2}$. It is
a labeling convention rather than a physical input --- fixing the
normalization of the O($\lambda^{2}$) expansion before $b_{12}$ + $b_{13}$ is determined.

\textbf{Condition 2 (cubic-group structural relation).} Take

\textbf{$b_{23}$ = (4/3)($b_{12}$ + $b_{13}$)}

which with Condition 1 gives $b_{12}$ + $b_{13}$ = 3/4. Unlike Condition 1, this
is a substantive physical input: it relates the dominant U-odd
coefficient of $A_{2}$ to the U-even combination through a fixed numerical
factor 4/3.

The geometric interpretation is that

\textbf{4/3 = $2A^{2}$ = 2 $\sin^{2}$($\angle$(e\_i, $\hat{h}$))}

where A = $\sqrt{2/3}$ is the off-democratic projection of any generation axis
(Sec. 2.2). Note that $A^{2}$ = (d$-$1)/d = $C_{2}(T_{1})/d$ is a structural identity
for the vector representation of any rotation group --- for d = 3 this
gives $A^{2}$ = 2/3 from both the off-democratic projection (1 $-$ 1/d) and the
Casimir-over-dimension ratio ($C_{2}/d$). The factor 4/3 in Condition 2
admits a \emph{partner-count $\times$ per-partner projection} reading: the
U-even combination $E_{12}$+$E_{13}$ couples generation $e_{1}$ to its two partners
\{$e_{2}$, $e_{3}$\}, with the per-partner projection nominally $A^{2}$. However, the
load-bearing claim of this argument is the \emph{operator alignment} ---
that the substrate's second-order Yukawa structure has a specific $T_{2}$
direction matching the partner-count $\times$ projection picture. Operator
counting on $T_{1}\otimes T_{1}\to T_{2}$ in the supporting framework {[}18{]}
shows that this $T_{2}$ direction is determined by the substrate's specific
Yukawa structure rather than forced by representation theory: the
natural Higgs-democratic spurion does not produce Cond 2. We therefore
classify Cond 2 as currently solution-specific (Layer 2(b) in the
four-layer framing of {[}18, §8.3{]}) --- the empirical match
indicates the framework's bijection has the alignment property, but a
Layer 2(a) promotion would require explicit substrate-level Yukawa
derivation.

In standard effective field theory, the U-even and dominant U-odd
coefficients $b_{12}$+$b_{13}$ and $b_{23}$ would be independent Wilson coefficients.
The structural relation $b_{23}$ = (4/3)($b_{12}$+$b_{13}$) is more naturally
accommodated in frameworks where quantum mechanics is itself an emergent
phenomenon arising from a deeper deterministic structure {[}15, 16, 17{]}, since EFT coefficients can then be
related by substrate-level mechanisms even absent a protecting EFT
symmetry. The factorization 4/3 = 2 $\times A^{2}$ is consistent with this
picture; whether it follows rigorously from the cubic-lattice
realization of {[}18{]} is the open task referenced above.

We can sharpen the open task. In the AF basis (charged leptons diagonal
at LO), second-order perturbative diagonalization of M\_$\ell$ M\_$\ell$$^{\dagger}$
generates contributions to the $A_{2}$ coefficients from cross-terms
involving $A_{1}$. Direct computation shows that the framework's choice $A_{1}$ =
($\sqrt{2}/3$)($E_{12}-E_{13}$) produces a hierarchy-independent cross-term
contribution $b_{23}$\^{}cross = $-$1/9 in the m\_$e^{2}\to$ 0 limit, while $b_{12}$ and
$b_{13}$ receive no analogous cross-term contribution. Cond. 2 at the
observable level then translates to a relation on the substrate's direct
second-order corrections:

\textbf{$b_{23}$\^{}sub = (4/3)($b_{12}$\^{}sub + $b_{13}$\^{}sub) + 1/9}

The +1/9 augmentation is forced by $A_{1}'s$ perturbative structure; deriving
the substrate corrections that satisfy this from the cubic-lattice
Yukawa is the quantitative open task.

We treat Condition 2 as a substrate-level relation and quote our angle
predictions on the understanding that two of them ($\sin^{2}\theta_{12}$ and $\sin^{2}\theta_{23}$)
sit at the layered level: their structural form is forced by cubic-group
representation theory (Layer 1), but the specific coefficient values
depend on a Layer 2(b) substrate alignment. Empirical agreement at 0.$07\sigma$
on $\sin^{2}\theta_{12}$ is, in this sense, a nontrivial test of the bijection's
alignment property.

At the angle level, this is equivalent to

\textbf{$2\Delta_{12}$ + $\Delta_{23}$ = 0}

which is testable directly against data (Sec. 4.3). The combination $b_{12}-b_{13}$ remains free; we identify it in Sec. 4.2 with $\delta$\_CP (the natural
identification, since it is the only remaining O($\lambda^{2}$) free parameter that
does not enter any of the three angles). The present analysis therefore
predicts all three angles but not $\delta$\_CP.

\hypertarget{numerical-predictions}{%
\subsubsection{3.5 Numerical predictions}\label{numerical-predictions}}

Substituting $b_{23}$ = 1, $b_{12}$ + $b_{13}$ = 3/4, and $\lambda^{2}$ = 1/($2\pi^{2}$):

\begin{itemize}
\tightlist
\item
  \textbf{$\sin^{2}\theta_{12}$ = 1/3 $-$ 1/($4\pi^{2}$) = 0.30798}
\item
  \textbf{$\sin^{2}\theta_{13}$ = 4/($18\pi^{2}$) = 0.02252}
\item
  \textbf{$\sin^{2}\theta_{23}$ = 1/2 + 1/($2\pi^{2}$) = 0.55066}
\end{itemize}

The first formula for $\sin^{2}\theta_{13}$ follows from $A^{2}$ = 2/3 and $\lambda^{2}$ = 1/($2\pi^{2}$)
together with the constraints on $A_{1}$; it does not use the normalization
conditions on $A_{2}$. The other two additionally use Conditions 1 and 2.

\hypertarget{the-sinuxb2ux3b8ux2081ux2082-prediction}{%
\subsection{4. The $\sin^{2}\theta_{12}$
prediction}\label{the-sinuxb2ux3b8ux2081ux2082-prediction}}

\hypertarget{numerical-value-and-comparison-with-data}{%
\subsubsection{4.1 Numerical value and comparison with
data}\label{numerical-value-and-comparison-with-data}}

The headline prediction is

\begin{quote}
\textbf{$\sin^{2}\theta_{12}$ = 1/3 $-$ 1/($4\pi^{2}$) = 0.30798}
\end{quote}

Comparison with current measurements:

\begin{longtable}[]{@{}lll@{}}
\toprule\noalign{}
Source & $\sin^{2}\theta_{12}$ & $\sigma$ from prediction \\
\midrule\noalign{}
\endhead
\bottomrule\noalign{}
\endlastfoot
This work & 0.30798 & --- \\
JUNO direct & 0.3092 ± 0.0087 & 0.$14\sigma$ \\
Post-JUNO global & 0.3085 ± 0.0073 & 0.$07\sigma$ \\
\end{longtable}

The prediction agrees with the post-JUNO global fit at 0.$07\sigma$. Both
inputs entering the prediction are independently motivated: 1/3 is the
$A_{4}$ TBM zeroth-order value, and $\lambda^{2}$ = 1/($2\pi^{2}$) is the squared Cabibbo
angle, fixed independently by the cubic Brillouin geometry (matching the
observed $\lambda$ = 0.22508 to 0.04\%).

\hypertarget{comparison-with-column-preservation-patterns}{%
\subsubsection{4.2 Comparison with column-preservation
patterns}\label{comparison-with-column-preservation-patterns}}

Tribimaximal mixing {[}1{]} was the leading phenomenological
description of neutrino oscillation data until Daya Bay measured a
nonzero $\theta_{13}$ {[}2{]}. Since then, modified TBM patterns have
been the subject of a substantial literature {[}3, 4, 5, 6{]}. The two best-known modifications
preserve one column of the original TBM matrix exactly: TM1 (first
column) and TM2 (second column). Each gives a one-parameter family that
accommodates the observed $\theta_{13}$ but predicts different values for the
other two angles. Most existing TBM-modification proposals start from a
discrete \emph{flavor symmetry} --- typically $A_{4}$, $S_{4}$, or a $Z_{2}$ residual
--- postulated as an axiom of the lepton sector. The minimal
modifications then preserve one column of the TBM matrix or one residual
symmetry of the neutrino mass matrix, yielding one-parameter
correlations between angles. Post-JUNO, He {[}9{]} finds that TM2
is excluded at \textgreater3.$5\sigma$ while TM1 remains viable; Zhang {[}10{]} reports similar results.

Comparison summary:

\begin{longtable}[]{@{}
  >{\raggedright\arraybackslash}p{(\columnwidth - 6\tabcolsep) * \real{0.2632}}
  >{\raggedright\arraybackslash}p{(\columnwidth - 6\tabcolsep) * \real{0.4211}}
  >{\raggedright\arraybackslash}p{(\columnwidth - 6\tabcolsep) * \real{0.2368}}
  >{\raggedright\arraybackslash}p{(\columnwidth - 6\tabcolsep) * \real{0.0789}}@{}}
\toprule\noalign{}
\begin{minipage}[b]{\linewidth}\raggedright
Approach
\end{minipage} & \begin{minipage}[b]{\linewidth}\raggedright
Symmetry origin
\end{minipage} & \begin{minipage}[b]{\linewidth}\raggedright
$\sin^{2}\theta_{12}$
\end{minipage} & \begin{minipage}[b]{\linewidth}\raggedright
$\sigma$
\end{minipage} \\
\midrule\noalign{}
\endhead
\bottomrule\noalign{}
\endlastfoot
TM1 & postulated $A_{4}$ flavor & 0.3184 & 1.$35\sigma$ \\
TM2 & postulated $A_{4}$ flavor & 0.3408 & 4.$4\sigma$ \\
TM3 & postulated; $\theta_{13}$ = 0 & --- & excluded \\
\textbf{This work} & \textbf{cubic group O of 3D space} &
\textbf{0.3080} & \textbf{0.$07\sigma$} \\
\end{longtable}

(Comparing against post-JUNO global fit 0.3085 ± 0.0073. TM1 uses
$\sin^{2}\theta_{12}$ = 1 $-$ (2/3)/$\cos^{2}\theta_{13}$; TM2 uses $\sin^{2}\theta_{12}$ = 1/(3$\cos^{2}\theta_{13}$); both with
$\sin^{2}\theta_{13}$ = 0.02195.)

Three differences distinguish the present prediction from TM\_n:

\textbf{(i) Symmetry origin.} TM1, TM2, and TM3 take a discrete flavor
symmetry as an axiom of the lepton sector. The cubic-group prediction
here uses $A_{4}\subset$ O, where O is the symmetry group of three-dimensional
space and the three generations transform as the $T_{1}$ irrep (Sec. 2.1;
forced by the framework's Bravais-lattice uniqueness and
anomaly-cancellation chain).

\textbf{(ii) Independence of predictions.} TM1 gives a one-parameter
correlation, $\sin^{2}\theta_{12}$ = 1 $-$ (2/3)/$\cos^{2}\theta_{13}$, fixing $\theta_{12}$ given $\theta_{13}$ but
predicting neither separately. The cubic-group analysis predicts all
three angles as fixed numerical values, with no parameter common to them
(Sec. 3.5).

\textbf{(iii) CP violation.} He's surviving pattern predicts sin $\delta$\_CP =
±0.998, close to maximal. Here, $\delta$\_CP is identified with the only
remaining free parameter at O($\lambda^{2}$), namely $b_{12}-b_{13}$, which does not
enter any of the three angles. The angle predictions therefore stand
independently of any value of $\delta$\_CP, and the present paper makes no
prediction for it.

The bottom line: the cubic-group prediction is closer to the post-JUNO
global fit by a factor of about 20 in $\sigma$-distance than the best surviving
column-preservation pattern, and at the same time makes more falsifiable
claims (three independent angles plus the sum rule below).

\hypertarget{the-sum-rule-discriminator}{%
\subsubsection{4.3 The sum-rule
discriminator}\label{the-sum-rule-discriminator}}

The cubic-group prediction implies

\textbf{2$\sin^{2}\theta_{12}$ + $\sin^{2}\theta_{23}$ = 7/6 = 1.1667}

versus the observed 2(0.3085) + 0.561 = 1.178 ± 0.020 (using the
post-JUNO $\sin^{2}\theta_{12}$ and the NuFIT 6.0 NO best fit for $\sin^{2}\theta_{23}$). The
agreement is at 0.$6\sigma$. This sum rule is the angle-level statement of $b_{23}$
= (4/3)($b_{12}$ + $b_{13}$), which is structural in the cubic-group setup. TM1
and TM2 do not produce this combination as an exact sum rule; in TM2,
$\sin^{2}\theta_{12}$ is determined by $\theta_{13}$ alone with no $\theta_{23}$ correlation. A direct
measurement of 2$\sin^{2}\theta_{12}$ + $\sin^{2}\theta_{23}$ at the 0.005 level is therefore the
sharpest data-side test that distinguishes the two classes of TBM
modification.

\hypertarget{the-remaining-angles}{%
\subsection{5. The remaining angles}\label{the-remaining-angles}}

\hypertarget{sinuxb2ux3b8ux2081ux2083}{%
\subsubsection{5.1 $\sin^{2}\theta_{13}$}\label{sinuxb2ux3b8ux2081ux2083}}

The formula for $\sin^{2}\theta_{13}$ is fully derived: at O($\lambda^{2}$) the $A_{2}$ parameters
cancel from $|U_e3|^{2}$, leaving $\sin^{2}\theta_{13}$ = $A^{4}\lambda^{2}$ = 4/($18\pi^{2}$) =
0.02252. The NuFIT 6.0 best fit (normal ordering) is 0.02195 ± 0.00056,
agreeing at 1.$01\sigma$. Equivalently,

\textbf{$\sin\theta_{13}$ = $A^{2}\lambda$ = (2/3) $\cdot$ 1/($\pi\sqrt{2}$) = 0.1500}

giving $\theta_{13}$ = 8.63°, compared with the observed 8.56° ± 0.10°.

\hypertarget{sinuxb2ux3b8ux2082ux2083}{%
\subsubsection{5.2 $\sin^{2}\theta_{23}$}\label{sinuxb2ux3b8ux2082ux2083}}

The formula gives $\sin^{2}\theta_{23}$ = 1/2 + 1/($2\pi^{2}$) = 0.5507. NuFIT 6.0 (NO) gives
$\sin^{2}\theta_{23}$ = 0.561$^{+}$$^{0}\cdot ^{012}$$_{-}$$_{0}$.$_{015}$, agreeing at 0.$74\sigma$. The prediction lies in
the second octant, in line with the current global preference. The sum
rule 2$\sin^{2}\theta_{12}$ + $\sin^{2}\theta_{23}$ = 7/6 ties this directly to $\sin^{2}\theta_{12}$:

\textbf{$\sin^{2}\theta_{23}$ = 7/6 $-$ 2$\sin^{2}\theta_{12}$}

A ±0.0014 ultimate JUNO error on $\sin^{2}\theta_{12}$ then implies a ±0.0028 error on
the predicted $\sin^{2}\theta_{23}$, comparable to the current direct measurement.

\hypertarget{summary}{%
\subsubsection{5.3 Summary}\label{summary}}

\begin{longtable}[]{@{}llll@{}}
\toprule\noalign{}
Angle & Prediction & Measurement & $\sigma$ \\
\midrule\noalign{}
\endhead
\bottomrule\noalign{}
\endlastfoot
$\sin^{2}\theta_{12}$ & 1/3 $-$ 1/($4\pi^{2}$) = 0.3080 & 0.3085 ± 0.0073 & 0.07 \\
$\sin^{2}\theta_{13}$ & 4/($18\pi^{2}$) = 0.02252 & 0.02195 ± 0.00056 & 1.01 \\
$\sin^{2}\theta_{23}$ & 1/2 + 1/($2\pi^{2}$) = 0.5507 & 0.561$^{+}$$^{0}\cdot ^{012}$$_{-}$$_{0}$.$_{015}$ & 0.74 \\
\end{longtable}

All three predictions agree with observation within 1.$1\sigma$.

\hypertarget{robustness}{%
\subsection{6. Robustness}\label{robustness}}

\hypertarget{cumulative-empirical-match}{%
\subsubsection{6.1 Cumulative empirical
match}\label{cumulative-empirical-match}}

The 0.$07\sigma$ match for $\sin^{2}\theta_{12}$ does not stand alone. It sits within a set
of empirical matches all using the same cubic-group structural inputs:

\begin{longtable}[]{@{}
  >{\raggedright\arraybackslash}p{(\columnwidth - 8\tabcolsep) * \real{0.2000}}
  >{\raggedright\arraybackslash}p{(\columnwidth - 8\tabcolsep) * \real{0.2000}}
  >{\raggedright\arraybackslash}p{(\columnwidth - 8\tabcolsep) * \real{0.2000}}
  >{\raggedright\arraybackslash}p{(\columnwidth - 8\tabcolsep) * \real{0.2000}}
  >{\raggedright\arraybackslash}p{(\columnwidth - 8\tabcolsep) * \real{0.2000}}@{}}
\toprule\noalign{}
\begin{minipage}[b]{\linewidth}\raggedright
Observable
\end{minipage} & \begin{minipage}[b]{\linewidth}\raggedright
Structural input
\end{minipage} & \begin{minipage}[b]{\linewidth}\raggedright
Prediction
\end{minipage} & \begin{minipage}[b]{\linewidth}\raggedright
Observed
\end{minipage} & \begin{minipage}[b]{\linewidth}\raggedright
Match
\end{minipage} \\
\midrule\noalign{}
\endhead
\bottomrule\noalign{}
\endlastfoot
Cabibbo angle & Brillouin distance $\pi\sqrt{2}$; chirality 1/\textbar q\textbar{}
& $\lambda$ = 1/($\pi\sqrt{2}$) = 0.22508 & 0.22500 ± 0.00067 & 0.04\% \\
Koide angle (charged leptons) & $C_{2}/d^{2}$ for cubic group & $\theta_{0}$ = 2/9 =
0.22222 & 0.222204 ± 0.000010 & 0.02\% \\
Solar mixing $\sin^{2}\theta_{12}$ (this work) & $A_{4}$ representation 1/3 + Cond 2 cubic
relation & 1/3 $-$ 1/($4\pi^{2}$) = 0.30798 & 0.3085 ± 0.0073 & 0.$07\sigma$ \\
PMNS sum rule (this work) & Cond 2 cubic relation & 2$\sin^{2}\theta_{12}$ + $\sin^{2}\theta_{23}$ =
7/6 & 1.178 ± 0.020 & 0.$6\sigma$ \\
Reactor mixing $\sin^{2}\theta_{13}$ (this work) & $A^{2} \lambda$ projection geometry & 4/($18\pi^{2}$)
= 0.02252 & 0.02195 ± 0.00056 & 1.$01\sigma$ \\
Atmospheric mixing $\sin^{2}\theta_{23}$ (this work) & Cond 2 + sum rule & 1/2 +
1/($2\pi^{2}$) = 0.55066 & 0.561$^{+}$$^{0}\cdot ^{012}$$_{-}$$_{0}$.$_{015}$ & 0.$74\sigma$ \\
\end{longtable}

Two structural parameters appear repeatedly: $\lambda^{2}$ = 1/($2\pi^{2}$), the squared
Cabibbo angle from the cubic Brillouin geometry, and $A^{2}$ = 2/3, the
off-democratic projection $\sin^{2}$($\angle$(e\_i, $\hat{h}$)) from the cubic-group
representation. Both enter multiple observables at sub-percent or sub-$\sigma$
precision. Under any reasonable prior over flavor models, this is
non-trivial structural agreement: a single set of inputs produces six
distinct empirical matches, each of which would have to be explained as
coincidence in a flavor-symmetry framework that postulates the cubic
structure rather than deriving it.

\hypertarget{sensitivity-to-condition-2}{%
\subsubsection{6.2 Sensitivity to Condition
2}\label{sensitivity-to-condition-2}}

The factor 4/3 in Condition 2 (Sec. 3.4) is the one substantive
structural input. Its observable consequences enter the sum rule $2\Delta_{12}$ +
$\Delta_{23}$ = 0 and, through Conditions 1 and 2 jointly, the absolute
predictions for $\sin^{2}\theta_{12}$ and $\sin^{2}\theta_{23}$.

A small shift to 4/3 $\to$ (4/3)(1 + $\epsilon$) propagates to $\Delta_{12}\to\Delta_{12}$ +
(1/2)$\cdot\epsilon\cdot \Delta_{23}$ at first order in $\epsilon$. With $\Delta_{23}\approx$ 1/($2\pi^{2}$) $\approx$ 0.05, a fractional
shift $\epsilon$ = 0.10 (a 10\% shift in Cond 2's structural coefficient) gives a
shift in $\sin^{2}\theta_{12}$ of \textasciitilde0.0025 --- comparable to JUNO's
ultimate-precision uncertainty of ±0.0014 and roughly half of JUNO's
first-result uncertainty of ±0.0087. The current 0.$07\sigma$ agreement
therefore tests Cond 2 at the few-percent level on its structural
coefficient --- a non-trivial test, since 4/3 is precisely $2A^{2}$ where $A^{2}$
= 2/3 is the projection factor independently fixed by the cubic
representation geometry.

\hypertarget{status-of-condition-2}{%
\subsubsection{6.3 Status of Condition 2}\label{status-of-condition-2}}

Condition 2 is presented in Sec. 3.4 as a substrate-level relation
supported by the projection-geometry argument (4/3 = $2A^{2}$ with the
partner-count reading) and by the broader cubic-lattice framework of
{[}18{]}. The structural identity $A^{2}$ = (d$-$1)/d = $C_{2}(T_{1})/d$ for
vector representations gives the underlying algebraic content. However,
operator counting on $T_{1}\otimes T_{1}\to T_{2}$ in the supporting framework shows that
Condition 2 itself is not forced by representation theory: the
substrate's $T_{2}$ direction in second-order Yukawa structure is
bijection-specific, and the natural Higgs-democratic spurion does not
reproduce Condition 2. We therefore classify it as solution-specific
(Layer 2(b) in the §8.3 four-layer framing of {[}18{]}) ---
the empirical match at 0.$07\sigma$ on $\sin^{2}\theta_{12}$ indicates the framework's
bijection has the alignment property required by Condition 2, but a
Layer 2(a) promotion would require explicit substrate-level Yukawa
derivation that has not been performed. The angle predictions $\sin^{2}\theta_{12}$
and $\sin^{2}\theta_{23}$ therefore sit at the layered level (Layer 1 form + Layer
2(b) coefficient) rather than as unconditional structural predictions.
The reactor angle $\sin^{2}\theta_{13}$ is independent of Condition 2 (the $A_{2}$
coefficients cancel from $|U_e3|^{2}$) and remains Layer 1
unconditional structural.

\textbf{Positive support for Layer 2(b) classification.} Beyond operator
counting, an explicit substrate calculation in {[}18, §7.3{]}
confirms the classification: for the natural cubic-symmetric staggered
Yukawa structure on $T_{1}$ BZ corners, both the leading and subleading
channels produce equal off-diagonal Y$^{(}$$^{2}$$^{)}$ entries by cubic symmetry,
failing Condition 2 by a hierarchy factor \textasciitilde17. The
framework's bijection therefore has structure beyond the natural
cubic-symmetric form --- Condition 2's empirical match indicates the
bijection introduces direction-dependent corrections that the natural
form lacks. The 0.$07\sigma$ match is therefore a tighter test of the
bijection's specific $T_{2}$-channel alignment than generic Layer 2(b)
classification suggests.

\textbf{Cubic-group consistency test.} The PMNS sum rule via Condition 2
is one observable in a six-prediction simultaneous-match cluster derived
from the structural identity $A^{2}$ = $C_{2}/d$ on the d=3 cubic lattice. The
other five are Wolfenstein A (1.2\%), $\sin^{2}\theta_{13}$ (0.4\%), Koide $\theta_{0}$
(0.02\%), m\_u/m\_d (0.6\%), and \textbar V\_cb\textbar{} (0.$4\sigma$). All
six match observation at sub-percent or sub-$\sigma$ precision simultaneously
{[}18, §7.7{]}. This places the sum rule's 0.$6\sigma$ match in a
broader pattern of cubic-group structural consistency across the
framework.

\textbf{Discretization-determined Layer 2(b).} Condition 2 is
therefore not a free parameter to be fit, but a derivable consequence
of the lattice's discretization scheme. The bijection space in
{[}18{]} is constrained by linear dynamics, T-invariance,
P-indivisibility, and cycle ergodicity, leaving a discretization-scheme
freedom (field quantization, lattice size, timestep, boundary
conditions, rounding rule) rather than abstract permutation freedom.
Layer 2(b) coefficients including Condition 2's $T_{2}$-direction alignment
are \emph{discretization-determined}, not independent Wilson
coefficients. The 0.$07\sigma$ empirical match at $\sin^{2}\theta_{12}$ and 0.$6\sigma$ at the sum
rule are tests of this discretization scheme.

The sum rule 2$\sin^{2}\theta_{12}$ + $\sin^{2}\theta_{23}$ = 7/6 is the most directly testable
consequence of Condition 2; its empirical match at 0.$6\sigma$ tests the
bijection's alignment property at the few-percent level on the
structural coefficient. The discriminating power of the sum rule against
TM1, TM2, and TM3 column-preservation patterns (Sec. 4.2, 4.3) is
independent of whether Condition 2 is Layer 2(a) or 2(b); the empirical
match is what matters for the discrimination.

\hypertarget{falsifiability}{%
\subsection{7. Falsifiability}\label{falsifiability}}

The prediction is sharp: $\lambda^{2}$ = 1/($2\pi^{2}$) is fixed independently by the
Cabibbo angle, and the coefficient 1/($4\pi^{2}$) in $\Delta_{12}$ has no adjustable
input. Three experimental directions test it.

\textbf{Sub-percent precision on $\sin^{2}\theta_{12}$.} Over its 30-year design
lifetime, JUNO is projected to reach $\sin^{2}\theta_{12}$ = 0.3092 ± 0.0014, a 6$\times$
improvement over the current first result. At that precision the
prediction 0.30798 is tested at 0.$86\sigma$ (using the JUNO 59-day central
value) or 0.$36\sigma$ (using the post-JUNO global central value). An ultimate
JUNO $2\sigma$ disagreement with 0.30798 would falsify the prediction.

\textbf{The sum rule.} The prediction 2$\sin^{2}\theta_{12}$ + $\sin^{2}\theta_{23}$ = 7/6 is
sharper than either angle alone: it ties JUNO's reactor measurement to
the atmospheric measurements at DUNE, T2HK, and Hyper-Kamiokande. As
both improve, the sum rule will be tested at the 0.005 level --- a
sharper test than either side separately.

\textbf{$\delta$\_CP.} The free parameter $b_{12}-b_{13}$ of $A_{2}$ does not enter any
angle prediction, so the framework is consistent with any $\delta$\_CP.
Long-baseline measurements at DUNE and Hyper-Kamiokande will determine
this remaining parameter directly.

\hypertarget{conclusion}{%
\subsection{8. Conclusion}\label{conclusion}}

We have derived $\sin^{2}\theta_{12}$ = 1/3 $-$ 1/($4\pi^{2}$) = 0.3080 from cubic-group flavor
structure, matching the post-JUNO global fit at 0.$07\sigma$ with no parameter
fit to $\sin^{2}\theta_{12}$. The same construction predicts $\sin^{2}\theta_{13}$ and $\sin^{2}\theta_{23}$
within 1.$1\sigma$ of the NuFIT 6.0 {[}14{]} best fits, and forces the exact sum rule
2$\sin^{2}\theta_{12}$ + $\sin^{2}\theta_{23}$ = 7/6, which distinguishes the prediction from the
TM1, TM2, and TM3 column-preservation patterns. The reactor angle
$\sin^{2}\theta_{13}$ = 4/($18\pi^{2}$) is unconditional structural (Layer 1 in {[}18, §8.3{]}), based on the representation-theoretic identity $A^{2}$ =
(d$-$1)/d = $C_{2}(T_{1})/d$ for vector reps. The solar and atmospheric angle
predictions are layered: their structural form is forced by cubic-group
representation theory (Layer 1), but the specific coefficient values
depend on Condition 2 ($b_{23}$ = (4/3)($b_{12}$+$b_{13}$)), which is currently
classified Layer 2(b) --- solution-specific to the bijection ---
following operator counting in the supporting framework. The empirical
match at 0.$07\sigma$ on $\sin^{2}\theta_{12}$ and at 0.$6\sigma$ on the sum rule indicates the
framework's bijection has the alignment property required by Condition
2; the predictions are therefore tests of that bijection-specific
structure as much as of the underlying cubic-group framework. The
Cabibbo scale $\lambda^{2}$ = 1/($2\pi^{2}$) entering the prediction is set independently
by cubic Brillouin-zone geometry and matches the observed Cabibbo angle
to 0.04\%.

JUNO's projected long-term precision will test the prediction at
sub-percent level over the next decade. Combined with atmospheric
measurements at DUNE, T2HK, and Hyper-Kamiokande, this will sharply test
the sum rule. A confirmed result at JUNO's ultimate precision would
support cubic-group flavor structure as a viable origin of the lepton
mixing pattern.

\hypertarget{acknowledgments}{%
\subsection{Acknowledgments}\label{acknowledgments}}

The author thanks the JUNO Collaboration for the rapid public release
of their first oscillation results, and F. Capozzi, E. Lisi, F.
Marcone, A. Marrone, and A. Palazzo for the prompt post-JUNO global
fit that enabled the present comparison. Discussions with the broader
emergent-quantum-mechanics community have informed the framework's
positioning relative to existing approaches; specific intellectual
debts are recorded in the references and in the underlying framework
documentation {[}18{]}.

\hypertarget{references}{%
\subsection{References}\label{references}}

{[}1{]} P. F. Harrison, D. H. Perkins, W. G. Scott, ``Tri-bimaximal
mixing and the neutrino oscillation data,'' Phys. Lett. B 530, 167
(2002).

{[}2{]} F. P. An et al.~(Daya Bay), ``Observation of
electron-antineutrino disappearance at Daya Bay,'' Phys. Rev.~Lett. 108,
171803 (2012).

{[}3{]} S. F. King and C. Luhn, ``Neutrino mass and mixing with discrete
symmetry,'' Rep.~Prog. Phys. 76, 056201 (2013).

{[}4{]} X.-G. He and A. Zee, ``Minimal modification to tribimaximal
mixing,'' Phys. Rev.~D 84, 053004 (2011).

{[}5{]} X.-G. He, ``TM1 mixing pattern from $A_{4}$,'' Chin. J. Phys. 53,
100101 (2015).

{[}6{]} S. T. Petcov and A. V. Titov, ``Assessing the viability of $A_{4}$,
$S_{4}$ and $A_{5}$ flavour symmetries for description of neutrino mixing,'' Phys.
Rev.~D 97, 115045 (2018).

{[}7{]} JUNO Collaboration (A. Abusleme et al.), ``First measurement of
reactor neutrino oscillations at JUNO,'' arXiv:2511.14593 (2025).

{[}8{]} F. Capozzi, E. Lisi, F. Marcone, A. Marrone, A. Palazzo,
``Updated bounds on the (1,2) neutrino oscillation parameters after
first JUNO results,'' arXiv:2511.21650 (2025).

{[}9{]} X.-G. He, ``Modified Tri-bimaximal neutrino mixing confronted by
JUNO $\theta_{12}$ measurement,'' arXiv:2511.15978 (2025).

{[}10{]} D. Zhang, ``Trimaximal mixing patterns meet the first JUNO
result,'' arXiv:2511.15654 (2025).

{[}11{]} E. Ma and G. Rajasekaran, ``Softly broken $A_{4}$ symmetry for
nearly degenerate neutrino masses,'' Phys. Rev.~D 64, 113012 (2001).

{[}12{]} G. Altarelli and F. Feruglio, ``Discrete flavor symmetries and
models of neutrino mixing,'' Rev.~Mod. Phys. 82, 2701 (2010).

{[}13{]} S. Navas et al.~(Particle Data Group), Phys. Rev.~D 110, 030001
(2024).

{[}14{]} I. Esteban, M. C. Gonzalez-Garcia, M. Maltoni, T. Schwetz, A.
Zhou, ``NuFIT 6.0: three-flavour global analyses of neutrino oscillation
experiments,'' JHEP (2024); www.nu-fit.org.

{[}15{]} G. 't Hooft, \emph{The Cellular Automaton Interpretation of
Quantum Mechanics}, Fundamental Theories of Physics, vol.~185 (Springer,
2016); arXiv:1405.1548.

{[}16{]} S. L. Adler, \emph{Quantum Theory as an Emergent Phenomenon:
The Statistical Mechanics of Matrix Models as the Precursor of Quantum
Field Theory} (Cambridge University Press, 2004).

{[}17{]} J. A. Barandes, ``The Stochastic-Quantum Correspondence,''
arXiv:2302.10778 (2023).

{[}18{]} A. Maybaum, \emph{The Incompleteness of Observation} (Zenodo
archive, 2026), https://doi.org/10.5281/zenodo.19060318.

{[}19{]} V. Chandrasekaran, R. Longo, G. Penington, E. Witten, ``An
algebra of observables for de Sitter space,'' \emph{JHEP} \textbf{02}
(2023) 082; arXiv:2206.10780.

{[}20{]} J. Maldacena, ``Real observers solving imaginary problems,''
arXiv:2412.14014 (2024).

{[}21{]} D. Harlow, M. Usatyuk, Y. Zhao, ``Quantum mechanics and
observers for gravity in a closed universe,'' \emph{JHEP} \textbf{02}
(2026) 108; arXiv:2501.02359.

{[}22{]} K. Slagle and J. Preskill, ``Emergent Quantum Mechanics at the
Boundary of a Local Classical Lattice Model,'' \emph{Phys. Rev.~A}
\textbf{108}, 012217 (2023); arXiv:2207.09465.

\end{document}
